\section{Abandonment: The Career of a Word}

``Abandonment'' is not a scriptural term and not a scholastic one. It is a word of French spiritual prose, and it acquired its technical sense in a single century---roughly between the publication of the \emph{Treatise on the Love of God} in 1616 and the condemnation of Fénelon's \emph{Maximes} in 1699. Everything afterwards, including the books an English-speaking Catholic is likely to be handed today, descends from that century and carries its scars. The word therefore has to be read historically before it can be read doctrinally.

\subsection{Francis de Sales: Two Wills and a Ball of Wax}

The Salesian account is the one the later tradition builds on, and its decisive move is a distinction between two ways God's will meets us. There is the \emph{signified will}---what God has made known that he wants believed, hoped, feared, loved, commanded, and counselled---and the \emph{will of good-pleasure}, which is God's absolute will, known to us chiefly by events. Books~VIII and IX of the \emph{Treatise} take them in turn: Book~VIII on conformity to the signified will, Book~IX on submission to the will of good-pleasure.\endnote{Francis de Sales, \emph{Treatise on the Love of God}, in the public-domain English translation of Henry Benedict Mackey, O.S.B.; text read 2026-07-25 at the Christian Classics Ethereal Library presentation (\texttt{ccel.org}, \emph{Treatise on the Love of God}, complete text). Book and chapter loci are given as printed there; the translator's own summary notes that Book~VIII ``treats of obedience to the already signified will of God'' and is completed by Book~IX, which ``treats of indifference, or the state of perfect readiness to accept all that God's good-pleasure may choose to send us.''}

Two features of the account do more work than their gentle tone suggests. The first is that the signified will proceeds ``by way of desire, and not by way of absolute will,'' so that we can obey or resist it. De Sales analyses this into three acts: ``he wills that we should be able to resist, he desires that we should not resist, and yet allows us to resist if we please.'' When we resist, ``God contributes nothing to our disobedience, but leaving our will in the hands of its liberty permits it to make choice of evil''; when we obey, ``God contributes his assistance, his inspiration, and his grace.'' Permission itself is ``an action of the will which of itself is barren, sterile and fruitless, and is as it were a passive action, which acts not but only permits action.'' And the divine desire is a real desire: God ``solicits, exhorts, excites, inspires, aids and succours us''---for one who forced the point would be treating his friend not as a guest but ``like a capon that has to be fattened.''\endnote{\emph{Treatise}, Book~VIII, ch.~III, as in the Mackey translation cited above; wording read 2026-07-25. The banquet image belongs to the same passage: a host who ``should by main force open his friend's mouth, cram meat into his throat, and make him swallow it, would not be giving courteous entertainment to his friend.''}

The second is that surrender to the will of good-pleasure is bounded, in de Sales's own formulation, by an exception. ``Nothing, except sin, is done without that will of God which is called absolute, or will of good-pleasure, which no one can hinder, and which is only known to us by events.''\endnote{\emph{Treatise}, Book~IX, ch.~I; read 2026-07-25. The clause ``except sin'' is the same exception the seventeenth-century devotional prose later abbreviates---and the same one modern readers of that prose most often lose.} That single clause is the hinge on which the whole doctrine turns, and it is the clause the later devotional literature abbreviates and readers of that literature drop.

The heart of Book~IX is chapter~IV, on indifference, and it opens by refusing to let indifference be confused with resignation. ``Resignation prefers God's will before all things, yet it loves many other things besides the will of God. Indifference goes beyond resignation: for it loves nothing except for the love of God's will.'' Then the image that made the doctrine famous: ``The indifferent heart is as a ball of wax in the hands of its God, receiving with equal readiness all the impressions of the Divine pleasure; it is a heart without choice, equally disposed for everything, having no other object of its will than the will of its God.'' And a page later comes the sentence that would cause a century of trouble: the indifferent soul ``would prize hell more with God's will than heaven without it; nay she would even prefer hell before heaven if she perceived only a little more of God's good-pleasure in that than in this, so that if by supposition of an impossible thing she should know that her damnation would be more agreeable to God than her salvation, she would quit her salvation and run to her damnation.''\endnote{\emph{Treatise}, Book~IX, ch.~IV; read 2026-07-25. The rhetorical form matters: the sentence is governed by ``by supposition of an impossible thing''---\emph{par supposition d'une chose impossible}---which marks it as a counterfactual whose antecedent de Sales denies. Its later career, when the qualifier is dropped, is treated in the section on boundaries below.}

What keeps the ball of wax from being quietism is the chapter that follows. De Sales's indifference is not indifference to what one is to do; it is indifference in the face of what one has not been told. ``The divine good-pleasure is scarcely known otherwise than by events, and as long as it is unknown to us, we must keep as close as possible to the will of God which is already declared or signified to us.'' His own worked example is a sickbed: he does not know whether God intends his mother's illness to end in death, ``but I know well that while awaiting the event from his good-pleasure, he wills, by his declared will, that I use remedies proper to effect a cure.'' Only when the event has actually occurred is there anything to acquiesce in. He then poses the objection directly: are we then to care for nothing and abandon our affairs to the mercy of events? ``Pardon me, Theotimus, we are to omit nothing which is requisite to bring the work which God has put into our hands to a happy issue, yet upon condition that, if the event be contrary, we should lovingly and peaceably embrace it.''\endnote{\emph{Treatise}, Book~IX, ch.~VI; read 2026-07-25. The chapter adds the distinction on which the whole practical rule rests: ``we are commanded to have great care in what appertains to God's glory and to our charge, but we are not bound to, or responsible for, the event, because it is not in our power''---glossed by the parable's ``take care of him,'' which, as St.~Bernard observed, is not ``cure him.''} The Abraham passage in the same chapter makes the same point in narrative: the mark of the indifferent heart is that Abraham lowers the knife at the angel's word as readily as he raised it at God's, being ``equally ready to sacrifice or not to sacrifice his son.''

One further chapter deserves notice because it anticipates the worst misuse of surrender. Book~IX, chapter~VIII, treats how we unite our will with God's ``in the permission of sins''---and its counsel is the opposite of passivity. ``God sovereignly hates sin, and yet he most wisely permits it, in order to let the reasonable creature act according to the condition of its nature''; and therefore, ``however obstinate sinners may be, let us never desist from aiding and assisting them.'' Only after tears and the offices of charity have been spent are we to turn our minds to other work.\endnote{\emph{Treatise}, Book~IX, ch.~VIII; read 2026-07-25. The chapter is the Salesian answer to the question whether conformity to God's permission of another's sin means acquiescence in that sin: it does not, and the required response is intercession and active help, not equanimity.}

\subsection{Guyon, Bossuet, Fénelon: The Quarrel That Set the Vocabulary}

At the end of the same century the vocabulary went to court. Jeanne Guyon's writings on pure love and passive prayer drew Fénelon's theological defence and Bossuet's theological attack, and the resulting quarrel---conducted at Versailles as much as in the schools---settled what French Catholic spirituality could safely say for two centuries afterwards.

The questions in dispute were not trivial ones. What does union with God mean? Can a human being be so detached as to will only what God wills? Can God be loved with a wholly disinterested love? Is union the normal flowering of a committed spiritual life, as Fénelon held, or a favour granted to particular predestined souls, as Bossuet held?\endnote{The characterisation of the disputed questions, and of Guyon's own sense of standing in a tradition running from the Rhineland mystics and John of the Cross through Francis de Sales, follows Dominique Salin, S.J., ``The Treatise on Abandonment to Divine Providence,'' \emph{The Way} 46/2 (April 2007): 21--36, at 24--26; read 2026-07-25 in the journal's freely posted PDF (\texttt{theway.org.uk}). The article is in copyright; it is used here as a named scholarly witness and quoted only in short phrases. Salin notes that ``abandonment'' was for Guyon and Fénelon ``a French culmination of a tradition that had started in Greek, passed through Latin and Flemish.''} Guyon took herself to be describing what earlier writers had called receptivity to the divine, the \emph{Gelassenheit} of the Rhinelanders, Ignatian indifference, or Salesian abandonment. Whether she was right about that is precisely what Rome had to decide.

It decided twice, and differently, in ways the next section examines proposition by proposition. Molinos was condemned in 1687 as a heretic and imprisoned; Fénelon's \emph{Explication des maximes des saints} was condemned in 1699 by a brief that carefully avoided calling him one, and he submitted at once from his own pulpit. The effect, as one modern historian of the literature puts it, was that ``for some two centuries subsequently, `mystical' language and themes vanished from spiritual literature''; writers spoke of devotion and piety rather than mysticism, and even the great Spanish mystics were handled with tongs.\endnote{Salin, ``The Treatise on Abandonment,'' 26, at the witness cited above. The same passage records that Guyon's ``more daring affirmations, though condemned within the Catholic Church, found their way into German Pietism and English Methodism.''} That reticence is the immediate literary background of the books this article is about.

\subsection{Caussade: A Classic Without a Settled Author}

The best-known modern treatise of abandonment is \emph{L'Abandon à la providence divine}, published in 1861 by the Jesuit Henri Ramière and attributed by him to Jean-Pierre de Caussade, S.J.\ (1675--1751). Its attribution history is a cautionary tale in the same genre as the anthology examined earlier in these pages, and it is now well documented.

The manuscript Ramière worked from had survived the Revolution and was kept by the Sisters of Nazareth at Montmirail, having originally belonged to a Paris Visitation convent. A note described it as letters ``written by an ecclesiastical person to a superior of a religious community,'' and someone had added that the author was Father Caussade of the Society of Jesus. Ramière took the indications at face value. He also supplied the title: the manuscript's own name, \emph{A Treatise Where One Discovers the True Art of the Perfection of Salvation}, ``gives no idea,'' he wrote, ``of the subject that it deals with.'' He reordered the material to some extent, with an eye to protecting unschooled readers from ideas that were too mystical or risky, and he noted in his preface that the treatise does not sufficiently stress that God's sanctifying action never occurs apart from the person of Jesus Christ. Later reprints added letters published as Caussade's, at least trebling the volume's size and creating an impression of a prolific author.\endnote{Salin, ``The Treatise on Abandonment,'' 21--23, at the witness cited above; the quoted phrases are Salin's rendering of Ramière's preface and of the manuscript note. Salin adds that a text closer to the original manuscript was published only in the 1960s, edited by Michel Olphe-Galliard, S.J.\ (1900--1985), who also separated the treatise from the two further volumes of letters, and that Salin himself produced a corrected edition in 2005.}

Modern research has gone further. Jacques Gagey's study of the manuscripts, and of the correspondence of a Nancy Visitation archivist, Sœur Fervel, who had been in contact with Ramière, shows that it was Fervel who persuaded Ramière that the treatise had been assembled from Caussade's letters to Mère Marie-Anne-Sophie de Rottembourg, and who passed him further manuscript letters as Caussade's. On Gagey's showing only thirty-two letters of spiritual direction are genuinely Caussade's, addressed to an unnamed ``lady from Lorraine''; and their unpolished style is, in Salin's judgement, ``far removed from the lyricism marking the treatise.'' Gagey attributes the treatise itself to that lady; Salin finds the arguments unconvincing and declines to name an author at all, observing that we know neither the writer's gender nor his or her ecclesiastical status, only that the writer had been steeped in the tradition of Francis de Sales and Madame Guyon.\endnote{Salin, ``The Treatise on Abandonment,'' 22--24, at the witness cited above, reporting Jacques Gagey, ed., \emph{L'Abandon à la providence divine d'une dame de Lorraine au XVIIIe siècle, suivi des Lettres spirituelles de Jean-Pierre Caussade à cette dame} (Grenoble: Jérôme Million, 2001). Gagey's volume and Olphe-Galliard's editions were not consulted directly for this article; they are cited as reported at the named witness, and no claim finer than what Salin reports is made here. Salin's own verdict is that ``it now seems almost impossible that the author was in fact the Jesuit, Jean-Pierre de Caussade.''}

The situation, then, is this. A classic of Catholic spirituality circulates under the name of a Jesuit whose authorship its most recent scholarly editor regards as almost impossible; in a text arranged and partly reordered by a nineteenth-century editor; under a title that editor invented; often bound with letters that mostly are not the named author's either. None of that makes the book bad. It makes it an anonymous eighteenth-century work of the Salesian--Guyonian tradition, transmitted through a nineteenth-century Jesuit's editorial judgement---which is what a careful reader should know before citing ``de Caussade'' as an authority for anything.

\subsection{Ramière's Fence}

What Ramière did next is more interesting than his attribution, and it is available in a public-domain English translation. To the 1861 edition he prefixed a preface headed ``Foundation and true nature of the virtue of abandonment, to explain and defend Father Caussade's doctrine''---in effect, a fence built around the treatise before the reader is let into it. Its opening sentences say why: ``There is no truth however clear which does not become error the moment it is lessened or exaggerated; and there is no food however salutary for the soul which may not, when ill-applied, become a fatal poison. The virtue of abandonment does not escape this danger.'' He then names the danger by name---the seventeenth-century Quietists, whose ``specious sophistries'' had for a time misled even ``the pious Fénelon''---and undertakes to set out ``the just limits which should mark our abandonment to divine Providence.''\endnote{\emph{Abandonment; or, Absolute Surrender to Divine Providence. Posthumous Work of Rev.~J.~P. de Caussade, S.J., Revised and Corrected by Rev.~H. Ramière, S.J.}, translated from the eighth French edition by Ella McMahon (New York, Cincinnati and Chicago: Benziger Brothers, 1887), preface, pp.~3--4; public-domain scan and OCR read 2026-07-25 at the Internet Archive, item \texttt{abandonmentorabs00caus} (imprimatur of Archbishop Michael Augustine, New York, 15 February 1887; the item's rights status is recorded as not in copyright). All Ramière quotations in this subsection are from that preface.}

The fence has three principles and a set of rules. The principles: nothing is done or happens, in the material or moral world, that God has not foreseen from eternity and has not willed or at least permitted; God can will or permit nothing except in view of the end he proposed in creating, his glory and the glory of Christ; and, Ramière's own addition, God wills to be glorified through this creature's happiness, so that in God's designs man's sanctification and happiness are inseparable from the divine glory. From these he draws the conclusion that confidence ``cannot be too great, too absolute, too child-like''---and then, immediately, the qualification that matters: ``This confidence in the fatherly providence of God cannot, evidently, dispense us from doing all that is in our power to accomplish His designs; but after having done all that depends upon our efforts we will abandon ourselves completely to God for the rest.''

The rules distinguish classes of events, and the distinctions are exactly the ones a modern reader needs. Where an event depends solely on God and human liberty had no part in producing or averting it---weather, unforeseeable accidents, natural defects of body or soul---``it is evident that our abandonment cannot be too absolute.'' Where the suffering comes through the malice of others, four things are to be done and kept apart: we cannot approve the offence against God, and should detest it, ``not because it wounds our self-love, but because it is an offence against the divine rights''; we may regard what is in itself an evil as a blessing for ourselves, by looking past the creature who caused it to the God who permitted it; we resign ourselves as to consequences already consummated and irreparable; and---this is the fourth rule, and the decisive one---``if it is in our power to avert the consequences of malice and injustice, and if in our true interest, and in the interest of the divine glory, we deem it necessary to take any measures to this end, let us do so without departing from the practice of the holy virtue of abandonment.'' Where our own faults are concerned, Ramière insists on distinguishing in the fault ``what is culpable and what is humiliating'': the culpability is to be hated and corrected, the humiliation accepted.

Finally he takes up the question that had destroyed Fénelon: should abandonment extend to accepting one's own damnation, and so make God the absolute sacrifice of every personal interest? ``To this point would Fenelon have carried the purity of love''---and, Ramière answers, ``the Church has condemned this doctrine which, in proposing to man a perfection contrary to his nature, reverses the order of God's designs.'' His argument is not merely juridical: ``How could one of the theological virtues be contrary to another, and charity exclude hope? What is eternal happiness if not the eternal reign of pure love?'' What perfect abandonment asks is that we observe in our desires the order of God's designs---glory first, our happiness inseparably second, never separated even hypothetically.

There is a certain irony in a nineteenth-century editor needing to erect this fence around a text he had himself misattributed, retitled, and rearranged. But the fence is sound, it is contemporaneous with the treatise's entry into modern devotional life, and it makes explicit what the tradition had learned the hard way. The next section examines what, exactly, it had learned.
